\documentclass[a4paper]{article}
% Español (México) con XeLaTeX: fontspec para fuentes Unicode, polyglossia
% para la división silábica y los nombres del documento en la variante mexicana.
\usepackage{fontspec}
\usepackage{polyglossia}
\setdefaultlanguage[variant=mexican]{spanish}
% Los nombres chinos van como «pinyin (original)»: xeCJK da a los caracteres
% Han su propia fuente; sin ella XeLaTeX los omite con un aviso.
\usepackage{xeCJK}
\setCJKmainfont{FandolSong-Regular.otf}
% polyglossia no traduce los nombres de listings.
\AtBeginDocument{\providecommand{\lstlistingname}{}\renewcommand{\lstlistingname}{Listado}%
  \providecommand{\lstlistlistingname}{}\renewcommand{\lstlistlistingname}{Índice de listados}}
\usepackage{amsmath, amssymb}
\usepackage{graphicx}
\usepackage[margin=2.5cm]{geometry}
\usepackage[most]{tcolorbox}
\usepackage{etoolbox}
\usepackage{listings}
\usepackage{hyperref}
\usepackage{booktabs}
\usepackage{float}
\newtcolorbox{knowledgebox}[1]{enhanced, colback=blue!5!white, colframe=blue!75!black, colbacktitle=blue!75!black, coltitle=white, fonttitle=\bfseries, title={#1}, attach boxed title to top left={yshift=-2mm, xshift=2mm}, boxrule=1pt, sharp corners}
\newtcolorbox{importantbox}[1]{enhanced, colback=yellow!10!white, colframe=yellow!80!black, colbacktitle=yellow!80!black, coltitle=black, fonttitle=\bfseries, title={#1}, sharp corners}
\newtcolorbox{warningbox}[1]{enhanced, colback=red!5!white, colframe=red!75!black, colbacktitle=red!75!black, coltitle=white, fonttitle=\bfseries, title={#1}, sharp corners}
\lstset{language=Python, basicstyle=\ttfamily\small, keywordstyle=\color{blue}, stringstyle=\color{red!60!black}, commentstyle=\color{green!60!black}, breaklines=true, frame=single, numbers=left, numberstyle=\tiny\color{gray}, captionpos=b, extendedchars=true}
\newcommand{\notetitle}{veRL en la práctica: entrenamiento SFT multi-turno}
\newcommand{\noteauthors}{Elaboradas a partir de materiales públicos del curso}
\newcommand{\notedate}{2025}
\newcommand{\videochannel}{Wudaokou Nash (五道口纳什)}
\newcommand{\videopublishdate}{2025}
\newcommand{\videoduration}{aproximadamente 20 minutos}
\newcommand{\videourl}{}
\newcommand{\videocoverpath}{cover.jpg}
\newcommand{\repourl}{https://github.com/hqhq1025/ai-course-notes}

% Footer with repo info
\usepackage{fancyhdr}
\pagestyle{fancy}
\fancyhf{}
\fancyfoot[C]{\thepage}
\fancyfoot[R]{\footnotesize\color{gray}\href{\repourl}{AI Course Notes}}
\renewcommand{\headrulewidth}{0pt}
\renewcommand{\footrulewidth}{0pt}

\begin{document}
\begin{titlepage}\centering
{\Large Notas del curso\par}\vspace{1.2cm}{\huge\bfseries \notetitle\par}\vspace{0.8cm}{\large \noteauthors\par}\vspace{0.3cm}{\large \notedate\par}\vspace{1.2cm}
\ifdefempty{\videocoverpath}{}{\includegraphics[width=0.82\textwidth,height=0.45\textheight,keepaspectratio]{\videocoverpath}\par}
\vfill
\begin{tcolorbox}[width=0.9\textwidth, colback=black!2!white, colframe=black!60, sharp corners]
\textbf{Autor o canal del video}: \videochannel\par\textbf{Fecha de publicación}: \videopublishdate\par\textbf{Duración del video}: \videoduration\par
\ifdefempty{\videourl}{}{\textbf{Enlace del video}: \href{\videourl}{\nolinkurl{\videourl}}\par}\par
\textbf{Página del proyecto}: \href{\repourl}{\nolinkurl{\repourl}}\par
\end{tcolorbox}
\end{titlepage}
\tableofcontents\newpage

\section{Introducción}
A partir de este episodio, se dedicarán varios episodios a presentar el proceso real de entrenamiento de veRL. En este episodio se presenta el entrenamiento Multi-turn SFT del proyecto ReTool de ByteDance.

\begin{importantbox}{Proyecto ReTool}
Trabajo de Coding Agent publicado por ByteDance en abril de 2025: el usuario plantea un problema de cómputo complejo y el modelo lo resuelve mediante un ciclo de escribir código, invocar el entorno, ejecutar el código y obtener retroalimentación. Se proporcionan el conjunto de datos completo, los scripts de entrenamiento y los pesos del modelo.
\end{importantbox}

\section{Tres ángulos de entrada}
Para analizar el entrenamiento de cualquier modelo, se aborda desde tres aspectos:
\begin{enumerate}
  \item \textbf{Data}: datos de Multi-turn tool-use
  \item \textbf{Algorithm}: SFT (fine-tuning supervisado)
  \item \textbf{Training}: configuración de entrenamiento de veRL
\end{enumerate}

\section{Características del SFT multi-turn}
Los datos de SFT contienen tool calls de múltiples turnos, cada turno con mensajes de user/assistant/tool. Durante el entrenamiento, la loss solo se calcula sobre las salidas del assistant.

\subsection{Resumen de la sección}
El SFT multi-turn es el paso previo del entrenamiento de Agentic RL, y aporta una política inicial adecuada para el entrenamiento de RL posterior.

\newpage
\section{Sugerencia de ingeniería: el SFT primero afianza las trayectorias de comportamiento}
En los escenarios de Agentic RL, el valor del SFT multiturno no es solo «entrenar primero un modelo que sepa hablar», sino afianzar primero las trayectorias de comportamiento básicas del uso de herramientas. Solo cuando el modelo ya ha aprendido a invocar herramientas correctamente en un diálogo de varios turnos, leer la retroalimentación y continuar generando, el RL posterior dejará de desperdiciar gran parte del presupuesto de muestreo en los errores de formato más elementales.

\begin{knowledgebox}{Por qué hacer primero SFT y después RL}
\begin{itemize}
  \item El SFT aprende primero la plantilla de comportamiento y reduce la dificultad de exploración del RL
  \item Los datos multiturno le enseñan al modelo a manejar la estructura temporal de los mensajes de herramienta
  \item Una buena política inicial eleva de manera notable la proporción efectiva de los rollouts posteriores
\end{itemize}
\end{knowledgebox}

\subsection{Resumen de la sección}
Para los sistemas agenticos, el SFT es el punto de partida de la alineación de comportamiento. Primero fija los patrones básicos de interacción, y después deja que el RL optimice la calidad de la política con un grano más fino.

\newpage
\section{Síntesis y ampliación}
\begin{enumerate}
  \item ReTool ofrece un caso completo de entrenamiento Agentic SFT $\to$ RL
  \item El procesamiento de datos multi-turn es un punto técnico clave
  \item Data + Algorithm + Training son los tres marcos para analizar el trabajo de entrenamiento
\end{enumerate}
\end{document}

